\documentclass[12pt, a4paper]{academics}

\academicsCoverLabel{解题报告}
\academicsCourse{算法设计与分析}{Design and Analysis of Algorithms}
\academicsReport{最短路径}{2026 年 6 月 5 日}
\academicsArthur{华博文}{19240212}

\begin{document}

\makeacademicscover

\section{引言}

最短路径问题是图论中最为经典、应用最为广泛的基础问题之一。其核心诉求看似简单——在图中寻找两顶点之间权值和最小的路径——但在不同约束条件下，这一问题衍生出丰富的算法变体，每种变体都推动着数据结构选择和建模方法论上的深化。

本报告以最短路径为主线，依次展开四道关联紧密的题目。从最基础的朴素 Dijkstra 算法在邻接矩阵上的 $\mathrm{O}(n^2)$ 贪心实现，到借助优先队列将复杂度优化至 $\mathrm{O}((n+m) \log n)$ 的堆优化版本；从引入负权边后 Dijkstra 贪心策略的失效，到 Bellman-Ford 的队列优化变体 SPFA 及其负环判定机制；最后，通过分层图的状态空间扩展技巧，解决带有限制条件的变种最短路问题。四道题目在数据范围上逐级扩大，在算法思想上层层递进：题一与题二展示数据结构选择对算法效率的根本性影响；题三探讨负权边带来的松弛次数隐患与负环判定的理论依据；题四引入分层图这一通用建模思想，将附加限制转化为状态空间的维度扩展。

报告中的所有代码均使用 C++23 标准编写，在 Apple clang 21.0.0 编译器下通过测试。

\section{题一：朴素 Dijkstra 最短路径}

\subsection{问题重述}

给定一个有向带权图，包含 $n$ 个顶点（编号 $1 \sim n$）和 $m$ 条有向边，所有边权均为非负整数。指定起点 $s$ 和终点 $t$，求从 $s$ 到 $t$ 的最短路径长度（即路径上各边权值之和的最小值）；若不可达，则输出 $-1$。

数据范围：$1 \le n \le 1000$，$1 \le m \le 10000$，边权不超过 $10^4$。输入格式：第一行为四个整数 $n, m, s, t$；接下来 $m$ 行，每行三个整数 $u, v, w$，表示一条从 $u$ 到 $v$、权值为 $w$ 的有向边。图中可能存在重边，处理时应取权值最小者。输出格式：一行一个整数，表示 $s$ 到 $t$ 的最短距离，若不可达则输出 $-1$。

\subsection{贪心选点与邻接矩阵松弛}

Dijkstra 算法的核心思想是贪心：在所有尚未确定最短距离的顶点中，选取当前 $\mathrm{dist}$ 值最小的顶点 $u$，将其标记为已确定，然后用 $u$ 的所有出边去松弛其邻居顶点。这一过程的正确性严格依赖于边权非负的前提——一旦某个顶点的 $\mathrm{dist}$ 值在当前所有未访问顶点中最小，就不可能有任何其他路径通过尚未访问的顶点以更短的代价到达它，因为所有中间边的权值均非负，任何绕道只会增加总代价而不会减少。

在 $n \le 1000$ 的数据规模下，朴素实现完全可行。算法共执行 $n$ 轮：每轮扫描全部 $n$ 个顶点以找出 $\mathrm{dist}$ 最小的未访问点（$\mathrm{O}(n)$），随后遍历所有 $n$ 个顶点进行松弛操作（同样 $\mathrm{O}(n)$）。使用邻接矩阵 $\mathit{g}[u][v]$ 存储边权，初始化为无穷大，读入时若遇到重边则取 $\min$ 以保证矩阵中保留的是最小权值。维护三个数组：$\mathrm{dist}[i]$ 记录当前已知的从 $s$ 到 $i$ 的最短距离（初始为无穷大，$\mathrm{dist}[s] = 0$）；$\mathrm{vis}[i]$ 标记顶点 $i$ 是否已确定最短距离；邻接矩阵 $\mathit{g}[u][v]$ 存储边权。

具体流程为：每轮循环先通过遍历 $\mathrm{vis}$ 为假且 $\mathrm{dist}$ 最小的顶点，将其标记为已访问。若该顶点的 $\mathrm{dist}$ 仍为无穷大，说明剩余顶点均不可达，可提前退出循环。随后遍历所有邻居进行松弛：若 $\mathrm{dist}[v] > \mathrm{dist}[u] + \mathit{g}[u][v]$，则更新 $\mathrm{dist}[v]$。最终 $\mathrm{dist}[t]$ 即为所求。

时间复杂度为 $\mathrm{O}(n^2)$——$n$ 轮选取，每轮扫描 $n$ 个顶点，松弛每条边一次。空间复杂度为 $\mathrm{O}(n^2)$，由邻接矩阵主导。当 $n = 1000$ 时，$n^2 = 10^6$ 次运算可在数十毫秒内完成。然而，当 $n$ 增长至 $10^5$ 时，$10^{10}$ 次运算将远远超出时限，这意味着我们必须转向稀疏图的专用数据结构——这正是下一道题目的核心课题。

\subsection{代码实现}

\inputminted{c++}{../src/p1_dijkstra_naive.cpp}

\section{题二：堆优化 Dijkstra 最短路径}

\subsection{问题重述}

本题的问题描述与题一完全一致——给定有向带权非负图，求 $s$ 到 $t$ 的最短路径长度——但数据范围显著扩大：$1 \le n \le 10^5$，$1 \le m \le 2 \times 10^5$，边权不超过 $10^4$。输入输出格式与题一相同。

\subsection{优先队列与惰性删除}

当 $n$ 达到 $10^5$ 时，题一的 $\mathrm{O}(n^2)$ 算法因邻接矩阵的空间开销和每轮扫描全部顶点的计算开销而双双失效。优化的方向是明确的：将数据结构从邻接矩阵升级为邻接表，将「查找最小 $\mathrm{dist}$ 顶点」的操作从线性扫描升级为优先队列。

邻接表将空间需求从 $\mathrm{O}(n^2)$ 压缩至 $\mathrm{O}(n+m)$，对于稀疏图（$m$ 与 $n$ 同阶）而言，这一升级是必要且充分的。而优先队列（最小堆）使得每次取出当前 $\mathrm{dist}$ 最小顶点的开销从 $\mathrm{O}(n)$ 降至 $\mathrm{O}(\log n)$，从根本上改变了算法的时间复杂度。

堆优化 Dijkstra 的流程为：将 $(0, s)$ 入堆；每次从堆顶取出距离最小的记录 $(d, u)$，若 $d \neq \mathrm{dist}[u]$，说明这是一条被后续更新淘汰的旧记录（即惰性记录），直接跳过；否则遍历 $u$ 的邻接表，对每条出边 $(u, v, w)$ 尝试松弛：若 $\mathrm{dist}[v] > \mathrm{dist}[u] + w$，则更新 $\mathrm{dist}[v]$ 并将 $(\mathrm{dist}[v], v)$ 入堆。

惰性删除是这一实现中的关键技巧。当一个顶点被多次松弛时，堆中可能保留多条以该顶点为目标的旧记录，每条都对应一个已被更新过的旧距离值。我们无需显式从堆中删除它们——仅凭 $d \neq \mathrm{dist}[u]$ 这一条件即可安全跳过过期记录，因为只有当前最优的距离值才有可能产生有效的松弛。这种技术简化了实现，且不增加渐进复杂度：每条边最多被松弛一次并产生一次入堆操作，堆的规模为 $\mathrm{O}(m)$，总时间复杂度为 $\mathrm{O}((n+m) \log n)$。对于 $n = 10^5, m = 2 \times 10^5$ 的规模，$(n+m) \log n \approx 3 \times 10^5 \times 17 \approx 5 \times 10^6$ 次运算完全可接受。空间复杂度为 $\mathrm{O}(n+m)$，由邻接表和堆共同决定。

\subsection{代码实现}

\inputminted{c++}{../src/p2_dijkstra_heap.cpp}

\section{题三：SPFA 与负权边最短路径}

\subsection{问题重述}

给定一个有向图，包含 $n$ 个顶点（$1 \le n \le 1000$）和 $m$ 条有向边（$1 \le m \le 5000$），边权 $w$ 可为任意整数（允许负数，但保证图中不存在从源点 $s$ 可达的负权环）。指定源点 $s$，计算 $s$ 到所有顶点的最短路径长度。若检测到从 $s$ 可达的负权环，则输出 \mil{NEGATIVE CYCLE}；否则，对于每个顶点按编号顺序输出其最短距离，不可达顶点输出 \mil{INF}（以空格分隔）。

输入格式：第一行为三个整数 $n, m, s$；接下来 $m$ 行，每行三个整数 $u, v, w$，表示一条从 $u$ 到 $v$、权值为 $w$ 的有向边。

\subsection{队列松弛与负环判定}

当图中存在负权边时，Dijkstra 算法的贪心基石不复存在。原因在于，一个已被标记为「已确定」的顶点的 $\mathrm{dist}$ 值，可能被后续经过负权边的路径进一步减小——负权边允许路径在绕道后反而获得更短的总距离。这一现象揭示了 Dijkstra 的适用边界：其正确性植根于「路径单调递增」的性质，而负权边的介入恰如其反。

Bellman-Ford 算法通过反复松弛所有边来应对负权边，其理论基础是：在一张不存在负环的图中，任意两点间的最短路径至多包含 $n-1$ 条边。因此，进行 $n-1$ 轮全边松弛后，所有最短路径必定确定；若第 $n$ 轮仍有顶点可被松弛，则必然存在负环。

SPFA（Shortest Path Faster Algorithm）是 Bellman-Ford 的队列优化版本，其核心改进在于：并非每轮无差别地扫描全部 $m$ 条边，而是仅将那些 $\mathrm{dist}$ 值刚被更新的顶点加入队列，由它们去松弛邻居。这一策略在随机稀疏图上的实际效率远优于标准 Bellman-Ford 的 $\mathrm{O}(nm)$，尽管最坏情况下复杂度仍为此量级。

负环判定的实现依赖于入队计数器 $\mathrm{cnt}[v]$：每当顶点 $v$ 松弛成功而被加入队列时，$\mathrm{cnt}[v]$ 自增；若某一顶点的入队次数超过 $n$，则表明存在负环。其原理是：在无负环图中，任意顶点的最短路径至多包含 $n-1$ 条边，因而每个顶点至多被松弛 $n-1$ 次；入队次数超过 $n$ 意味着最短路径的边数已突破理论上限，必然是由负环造成的无限松弛所致。

实现细节：维护 $\mathrm{dist}$ 数组（初始 $\infty$，$\mathrm{dist}[s]=0$）、$\mathrm{inq}$ 标志（记录顶点是否已在队列中，避免重复入队）、以及 $\mathrm{cnt}$ 计数器。每次从队列头部取出顶点 $u$，取消其 $\mathrm{inq}$ 标志，遍历其所有出边进行松弛；若松弛成功且邻居不在队列中，则将其入队并递增计数器，立即检查计数器是否超限。时间复杂度 $\mathrm{O}(nm)$，在 $n=1000, m=5000$ 的规模下约 $5 \times 10^6$ 至 $10^7$ 次运算；空间复杂度 $\mathrm{O}(n+m)$，由邻接表与辅助数组共同构成。

\subsection{代码实现}

\inputminted{c++}{../src/p3_spfa.cpp}

\section{题四：分层图最短路}

\subsection{问题重述}

给定一个无向带权连通图，包含 $n$ 个顶点（$1 \le n \le 10000$）和 $m$ 条无向边（$1 \le m \le 50000$），所有边权均为正整数。你有 $k$ 次机会（$0 \le k \le 10$）可以免费通过一条边（即不计该边的权值），每次机会只能用于一条边且不可拆分。求从起点 $s$ 到终点 $t$ 的最短路径总代价。

输入格式：第一行为五个整数 $n, m, k, s, t$；接下来 $m$ 行，每行三个整数 $u, v, w$，表示 $u$ 与 $v$ 之间有一条权值为 $w$ 的无向边。

\subsection{状态空间扩展与分层建模}

本题的核心挑战在于免费机会这一附加约束的引入。表面上看，一种自然的贪心思路是「直接免去路径上权值最大的 $k$ 条边」，但这一策略存在根本缺陷：被免费化的边必须恰好位于某条实际路径上，且免费机会的使用顺序与路径结构相互缠绕——后用的免费机会无法反过头去消减前面已走过的边的权值。贪心拆解无法保证最优子结构，因此必须寻找一种能够精确追踪免费机会使用状态的建模方法。

分层图正是为此类「带有附加状态维度」的最短路变体而生的通用技巧。其核心思想是将附加约束提升为状态空间的一个新维度：我们构造一个拥有 $(k+1) \times n$ 个状态的扩展图，其中状态 $(v, j)$ 表示「当前位于顶点 $v$，且已使用了 $j$ 次免费机会」。对于原图中的每条无向边 $(u, v, w)$，在分层图中产生两类转移：正常行走（付费），即从 $(u, j)$ 到 $(v, j)$，代价为 $w$；使用免费机会，即从 $(u, j)$ 到 $(v, j+1)$，代价为 $0$。转移是双向的（因为原边无向），且 $j$ 不能超过 $k$。

在这个扩展后的分层图上运行堆优化 Dijkstra 算法——所有边权仍然非负，因此 Dijkstra 的正确性得以保持——以 $(s, 0)$ 为起点。最终答案为 $\min_{0 \le j \le k} \mathrm{dist}[t][j]$，即到达终点时，无论使用了多少次免费机会（从 $0$ 到 $k$），所达到的最小总代价。

分层图建模的威力在于它的普适性：它将「带有 $k$ 次特殊权限」这一看似复杂的状态约束，通过状态空间的乘积扩展，整齐地映射回标准最短路框架。类似地，「必须经过某点集」「最多经过某条边」等约束均可通过状态扩展转化为标准最短路问题，这正是算法竞赛中「转化问题」这一元方法论的生动体现。

时间与空间复杂度均与分层图规模成正比。分层图共有 $(k+1)n$ 个状态和至多 $2(k+1)m$ 条边（每条原边对应两种转移，且无向边在两层间各出现一次）。堆优化 Dijkstra 在此图上的复杂度为 $\mathrm{O}((kn + km) \log (kn))$。由于 $k \le 10$ 为小常数，实际复杂度接近 $\mathrm{O}((n+m) \log n)$，对于 $n = 10000, m = 50000$ 的规模完全可行。空间方面，$\mathrm{dist}$ 二维数组占用 $\mathrm{O}(kn)$，邻接表 $\mathrm{O}(n+m)$，优先队列 $\mathrm{O}(km)$，总空间为 $\mathrm{O}(kn + m)$。

\subsection{代码实现}

\inputminted{c++}{../src/p4_layered_graph.cpp}

\section{总结}

本报告以最短路径为核心线索，选取了四道难度递进、思想关联的题目，系统展示了图论算法从基础实现到高阶建模的完整逻辑链条。

题一的朴素 Dijkstra 是所有后续讨论的起点。它以邻接矩阵为载体、以贪心选点为策略，直观地建立了单源最短路径的求解框架。$\mathrm{O}(n^2)$ 的复杂度在千级数据规模下绰绰有余，算法简明、正确性由边权非负的单调性直接保证，是初学者理解「松弛」与「贪心」两大核心概念的入门之作。

题二突破了朴素实现的瓶颈，通过将数据结构从邻接矩阵升级为邻接表加优先队列，将复杂度从 $\mathrm{O}(n^2)$ 压缩至 $\mathrm{O}((n+m) \log n)$。这一优化并非仅仅是常数级别的提升，而是量级上的跨越——它使得同一个 Dijkstra 算法能够处理从 $10^3$ 到 $10^5$ 级别的数据。惰性删除技术的引入进一步展示了工程实现中的巧思：以几乎为零的额外代码，安全地略过堆中的过期记录，不失渐进复杂度的同时保持了代码的简洁性。

题三将视野扩展至负权边场景。Dijkstra 的贪心假设在此失效，迫使我们回归 Bellman-Ford 的反复松弛框架。SPFA 作为其队列优化版本，通过仅维护松弛候选队列极大地减少了无效扫描，同时借助入队计数器实现了简洁而精确的负环判定。这一过程中，「最短路径边数不超过 $n-1$」这一理论上界作为负环判定的逻辑基础，启示我们：算法正确性的证明往往蕴含着实现优化的关键线索。

题四的分层图建模是本报告方法论上的制高点。面对「$k$ 次免费机会」这一附加约束，分层图技巧通过将免费次数提升为状态空间的第二维度，将原本需要复杂决策的问题整齐地映射回标准最短路框架。这一建模范式的威力在于其普适性——它揭示了一条处理带约束最短路问题的通用路径：不试图在原有算法中嵌入复杂的条件判断，而是扩展状态空间，在更大的图上运行标准算法。

回顾四道题目的递进脉络，可以看到两条贯穿始终的方法论线索。其一，数据结构选择决定算法效率的基底：从邻接矩阵到邻接表加堆，从线性扫描到对数查询，每一次数据结构升级都直接转化为复杂度量级的跃迁。其二，状态建模是处理复杂约束的核心手段：分层图将「免费机会数」提升为状态维度，正如动态规划中状态维度的增加往往伴随着问题表达力的增强——将附加约束编码为状态空间的维度扩展，是算法设计中一种具有高度可迁移性的思想范式。识别问题的结构约束，选择合适的数学模型与数据结构，在正确性与效率之间求得平衡，这正是算法设计与分析的核心要义。

\end{document}
